
%\documentclass[12pt,twocolumn]{amsart}
\documentclass[%
 reprint,
%superscriptaddress,
%groupedaddress,
%unsortedaddress,
%runinaddress,
%frontmatterverbose, 
%preprint,
%preprintnumbers,
%nofootinbib,
%nobibnotes,
%bibnotes,
 amsmath,amssymb,
 aps,
%pra,
%prb,
%rmp,
%prstab,
%prstper,
%floatfix,
]{revtex4-2}


% ~~~~~~~~~~~~~~~~~~~~~~~~~~~~~~~~~~~~~~~ V
%: -def.
\newcommand{\docName}{How to Study for CIS 111}
\newcommand{\AI}{\textsc{ai}}
\newcommand{\JDK}{\textsc{jdk}}
\newcommand{\IDE}{\textsc{ide}}
\newcommand{\cisOOO}{\textsc{cis111}}

% ~~~~~~~~~~~~~~~~~~~~~~~~~~~~~~~~~~~~~~~ A


\usepackage{graphicx}% Include figure files
\usepackage{dcolumn}% Align table columns on decimal point
\usepackage{bm}% bold math
%\usepackage{hyperref}% add hypertext capabilities
%\usepackage[mathlines]{lineno}% Enable numbering of text and display math
%\linenumbers\relax % Commence numbering lines

%\usepackage[showframe,%Uncomment any one of the following lines to test 
%%scale=0.7, marginratio={1:1, 2:3}, ignoreall,% default settings
%%text={7in,10in},centering,
%%margin=1.5in,
%%total={6.5in,8.75in}, top=1.2in, left=0.9in, includefoot,
%%height=10in,a5paper,hmargin={3cm,0.8in},
%]{geometry}



%% ~~~~~~~~~~~~~~~~~~~~~~~~~~~~~~~~~~~~~~~  V page layout
%%: -Pkg. page layout
%
%% ---------- Page layout ----------
%%\usepackage[margin=2cm]{geometry}
%\usepackage[
%  margin=2cm,
%  footskip=30pt,
%]{geometry}
%% ~~~~~~~~~~~~~~~~~~~~~~~~~~~~~~~~~~~~~~~ A page layout
%
%
%
%% ~~~~~~~~~~~~~~~~~~~~~~~~~~~~~~~~~~~~~~~ V encoding and fonts 
%%: -Pkg. encoding and fonts
%\usepackage[T1]{fontenc}
%\usepackage[utf8]{inputenc}
%\usepackage{lmodern}
%% ~~~~~~~~~~~~~~~~~~~~~~~~~~~~~~~~~~~~~~~ A encoding and fonts 
%
%
%% ~~~~~~~~~~~~~~~~~~~~~~~~~~~~~~~~~~~~~~~ V tables 
%%: -Pkg. table
%\usepackage{booktabs} % for much better looking tables
%\usepackage{tabularx}
%\usepackage{longtable}
%\usepackage{array} % for better arrays (eg matrices) in maths
%% ~~~~~~~~~~~~~~~~~~~~~~~~~~~~~~~~~~~~~~~ A tables 
%
%
% ~~~~~~~~~~~~~~~~~~~~~~~~~~~~~~~~~~~~~~~ V list 
%: -Pkg. list
\usepackage{enumitem}
% ~~~~~~~~~~~~~~~~~~~~~~~~~~~~~~~~~~~~~~~ A list 
%
%
%% ~~~~~~~~~~~~~~~~~~~~~~~~~~~~~~~~~~~~~~~ V links
%%: -Pkg. link
%\usepackage{hyperref}
%\hypersetup{
%  colorlinks=true,
%  urlcolor=blue,
%  linkcolor=blue,
%  citecolor=blue
%}
%% ~~~~~~~~~~~~~~~~~~~~~~~~~~~~~~~~~~~~~~~ A links



% ~~~~~~~~~~~~~~~~~~~~~~~~~~~~~~~~~~~~~~~ V listings
%: -Pkg. listing
%\usepackage{listings}
%\lstset{
%  basicstyle=\ttfamily,
%  keywordstyle=\color{blue},
%  language=Python,
%}

\usepackage{listings}
\usepackage{xcolor}

\lstset{
  basicstyle=\ttfamily,
  keywordstyle=\color{blue},
  commentstyle=\color{gray}\itshape,
  stringstyle=\color{red},
}

\newcommand{\code}[1]{\lstinline[language=Java]{#1}}
% ~~~~~~~~~~~~~~~~~~~~~~~~~~~~~~~~~~~~~~~ A listings



% ~~~~~~~~~~~~~~~~~~~~~~~~~~~~~~~~~~~~~~~ V timestamp
%: -Pkg. timestamp
%\usepackage[style=iso,showisoZ=true]{datetime2}    % yyyy-mm-dd \DTMnow
\usepackage[style=iso,showzone=false]{datetime2}    % yyyy-mm-dd \DTMnow
% ~~~~~~~~~~~~~~~~~~~~~~~~~~~~~~~~~~~~~~~ A timestamp



%% ~~~~~~~~~~~~~~~~~~~~~~~~~~~~~~~~~~~~~~~ V capitalized title
%%: -Pkg. capitalized title
%%\usepackage{sectsty}
%%\usepackage{titlecaps}
%%
%%\sectionfont{\normalfont\large\bfseries\Titlecap}
%%\subsectionfont{\normalfont\normalsize\bfseries\Titlecap}
%%\subsubsectionfont{\normalfont\normalsize\bfseries\Titlecap}
%
%
%%\usepackage{titlecaps}
%%
%%\makeatletter
%%\renewcommand\section{\@startsection{section}{1}%
%%  \z@
%%  {-1.5ex\@plus -.5ex \@minus -.2ex}%
%%  {.7ex \@plus .2ex}%
%%  {\normalfont\large\bfseries\Titlecap}}
%%
%%\renewcommand\subsection{\@startsection{subsection}{2}%
%%  \z@
%%  {-1.25ex\@plus -.5ex \@minus -.2ex}%
%%  {.5ex \@plus .2ex}%
%%  {\normalfont\normalsize\bfseries\Titlecap}}
%%
%%\renewcommand\subsubsection{\@startsection{subsubsection}{3}%
%%  \z@
%%  {-1.25ex\@plus -.5ex \@minus -.2ex}%
%%  {.5ex \@plus .2ex}%
%%  {\normalfont\normalsize\bfseries\Titlecap}}
%%\makeatother
%
%\makeatletter
%\renewcommand\section{\@startsection{section}{1}%
%  \z@
%  {-1.5ex\@plus -.5ex \@minus -.2ex}%
%  {.7ex \@plus .2ex}%
%  {\normalfont\large\bfseries}}
%
%\renewcommand\subsection{\@startsection{subsection}{2}%
%  \z@
%  {-1.25ex\@plus -.5ex \@minus -.2ex}%
%  {.5ex \@plus .2ex}%
%  {\normalfont\normalsize\bfseries}}
%
%\renewcommand\subsubsection{\@startsection{subsubsection}{3}%
%  \z@
%  {-1.25ex\@plus -.5ex \@minus -.2ex}%
%  {.5ex \@plus .2ex}%
%  {\normalfont\normalsize\bfseries}}
%\makeatother
%% ~~~~~~~~~~~~~~~~~~~~~~~~~~~~~~~~~~~~~~~ A capitalized title


% ~~~~~~~~~~~~~~~~~~~~~~~~~~~~~~~~~~~~~~~  V headers and footers
%: -Pkg. headers and footers
\newcommand{\headerfont}{\sffamily\small}
%
\usepackage{lastpage}   % <-- provides \pageref{LastPage}
%
\usepackage{fancyhdr}
\pagestyle{fancy}
\fancyhf{}
%\fancyhead[L]{aaa}
%\fancyhead[R]{aaa}
\fancyfoot[L]{\headerfont \docName}
\fancyfoot[C]{\headerfont \thepage\,/\,\pageref{LastPage}}
\fancyfoot[R]{\headerfont v\DTMnow}
%\lhead{}\chead{}\rhead{}
%\lfoot{}\cfoot{\thepage}\rfoot{}
\renewcommand{\headrulewidth}{0pt}
\renewcommand{\footrulewidth}{0.4pt}
% ~~~~~~~~~~~~~~~~~~~~~~~~~~~~~~~~~~~~~~~  A headers and footers

%% ---------- Paragraph spacing ----------
%\usepackage{parskip}
%% ~~~~~~~~~~~~~~~~~~~~~~~~~~~~~~~~~~~~~~~ A
%
%
%
%% See the ``Article customise'' template for come common customisations
%% +++++++++++++++++++++++++++++++++++++++ V
%%: -Fig.
%%% LaTeX - Article customise
%
%%%% PACKAGES
%\usepackage{verbatim} % adds environment for commenting out blocks of text & for better verbatim
%%\usepackage{subfigure} % make it possible to include more than one captioned figure/table in a single float
%\usepackage{subcaption} % modern






% ~~~~~~~~~~~~~~~~~~~~~~~~~~~~~~~~~~~~~~ V
\begin{document}

\title{How to Study for CIS 111}
\author{
	Haluk Bingol
}
\affiliation{
	Faculty of Computer and Information Sciences\\
	Yeditepe University
}
\date{\today}% It is always \today, today,
             %  but any date may be explicitly specified

% ~~~~~~~~~~~~~~~~~~~~~~~~~~~~~~~~~~~~~~~ V
\begin{abstract}
	Unfortunately K12 education is quite different then what is expected in university.
	K12 is focused on multiple choice questions.
	In university, you are expected to ``create'' new things from a ``blank'' paper.
	This is also the case in programming.
	You start with almost empty source file.
	
	First year students to the Faculty of Computer and Information Sciences suffer a lot from \cisOOO.
	\cisOOO\ seems to them very difficult because they are expected to design and implement a new application in the labs and in the exams,
	which are done on computer.
	This is an attempt to help our new students how to be successful in \cisOOO\ and, 
	in general, 
	in programming.	
\end{abstract}
% ~~~~~~~~~~~~~~~~~~~~~~~~~~~~~~~~~~~~~~~ A

\maketitle

\tableofcontents
% ~~~~~~~~~~~~~~~~~~~~~~~~~~~~~~~~~~~~~~ A








% ~~~~~~~~~~~~~~~~~~~~~~~~~~~~~~~~~~~~~~
\section{General Study Strategy}

Studying only the presentation slides is not enough for success in \cisOOO. 
Read the related chapter in the textbook. 
The book usually explains concepts more clearly, covers more material, and 
provides additional examples.

After reading, try to solve a few programming exercises at the end of the chapter. 
If you cannot solve them, go back to the relevant section and read it again. 
Then review the sample exam questions here:
\href{https://drive.google.com/drive/folders/1r1TcJdlwSAhWCe9vA9BqdJSkVTFXkmkM?usp=drive_link}
	{link}.









% ~~~~~~~~~~~~~~~~~~~~~~~~~~~~~~~~~~~~~~
\section{How to Solve a Programming Exercise}

\AI\ can sometimes solve an entire problem for you. 
If you choose that route, you will not learn how to program. 
You may memorize the solution, but that will not help you in exams.
To actually learn how to code the solution, follow this process:

\begin{enumerate}[label=\roman*)]
	
	\item 
	First, try to solve the problem by yourself.
	
	\item 
	If you cannot solve it, ask \AI\ for a hint about how to approach the problem. 
	Do not ask for the full code yet.
	
	\item 
	With that hint, try again by yourself.
	
	\item 
	If you still cannot solve it, go to the Help Desk. 
	Help Desk staff instructed not to solve the entire problem for you. 
	They will guide you toward the correct approach. 
	Bring your code, error messages, and notes about what you have tried.
	
	\item 
	Try again by yourself.
	
	\item 
	Only after these steps should you ask \AI\ for a complete solution. 
	When you receive it, study it, close it, and rewrite the solution from scratch without looking.
	
\end{enumerate}




% ~~~~~~~~~~~~~~~~~~~~~~~~~~~~~~~~~~~~~~
\section{Advice for First-Time Programmers}

A first-year student in their first programming course should focus on doing, not just reading.

\begin{itemize}
	
	\item
	\emph{Write code every day, even in small amounts.} 
	Programming is a skill, like playing an instrument. 
	You learn by practicing, not by watching.
	
	\item
	Read a concept, then immediately type out examples yourself. 
	Do not copy and paste.
	
	\item
	Type out examples from lectures and the textbook by hand. 
	Then modify them: change a value, break the code on purpose, and see what happens. 
	\emph{Understanding why something breaks teaches more than seeing it work.}
	
	\item
	\emph{Start assignments early, not the night before.} 
	Programming problems take unpredictable amounts of time, and bugs often need time to solve.
	
	\item
	\emph{Learn to read error messages instead of fearing them.} 
	They usually tell you the line and the problem. 
	Getting comfortable with debugging is half the job.
	
	\item
	\emph{Master the fundamentals deeply before moving on}: 
	variables, loops, conditionals, functions, and data types. 
	These appear in every language and later course.
	
	\item
	Explain concepts out loud or to a classmate. 
	If you cannot explain how a loop works, you do not fully understand it yet.
	
	\item
	Use office hours and do not struggle silently for more than about 30 minutes on a single bug. 
	\emph{Learning when to ask for help is a skill.}
	
	\item
	Practice beyond assigned work with small personal projects: 
	a calculator, a number guessing game, a to-do list. 
	\emph{Applying concepts to something you chose makes them stick.}
	
	\item
	The biggest mistake freshmen make is trying to memorize syntax or passively reading. 
	\emph{Programming rewards active, hands-on repetition.}

\end{itemize}




% ~~~~~~~~~~~~~~~~~~~~~~~~~~~~~~~~~~~~~~
\section{How to Study for a First Java Course}

	
\begin{itemize}

	\item
	\textbf{Set up your tools early.} 
	Install the \JDK\ and an \IDE\ such as IntelliJ IDEA Community or \textbf{Eclipse} in the first week. 
	Learn to compile and run from both the \IDE\ and the command line 
	(\code{javac} then \code{java}), 
	so you understand what is actually happening.
	
	\item
	\textbf{Get comfortable with Java’s structure right away.} 
	Every program lives inside a class, and execution starts at 
	\code{public static void main(String[] args)}. 
	This boilerplate confuses beginners because it contains keywords you do not understand yet. 
	Accept it early, then revisit each keyword (
	\code{public}, 
	\code{static}, 
	\code{void}
	) as the course explains them.
	
	\item
	\textbf{Master the fundamentals in order:} 
	variables and primitive types (
	\code{int}, 
	\code{double}, 
	\code{boolean}, 
	\code{char}
	), then 
	\code{String}, then operators, 
	conditionals (
	\code{if}/\code{else}, 
	\code{switch}), 
	loops (\code{for}, \code{while}), 
	arrays, and finally methods. 
	Do not move on until each feels automatic.
	
	\item
	\textbf{Pay close attention to types.} 
	Java is statically typed, so many beginner errors are type mismatches. 
	Understand the difference between integer and floating-point division: 
	\code{5 / 2} is 2, not 2.5. 
	Also understand why \code{==} on String values behaves unexpectedly and 
	why you should use 
	\code{.equals()} instead.
	
	\item
	\textbf{Understand objects and classes deeply.} 
	This is where Java courses separate students who pass from those who thrive. 
	Practice writing your own classes with fields, constructors, and methods. 
	Understand the difference between a class (the blueprint) and an object (the instance), and 
	what this refers to.
	
	\item
	\textbf{Type out and modify every example.} 
	Write a class, add a method, break it, read the compiler error, and fix it. 
	Java’s compiler errors are strict but informative. Learn to read them.
	
	\item
	\textbf{Practice with small programs}: 
	a temperature converter, 
	a grade calculator, 
	a simple bank account class, and 
	a number guessing game. 
	These exercise variables, loops, conditionals, and eventually objects.
	
	\item
	Common beginner traps to watch for:
	
	\begin{itemize}
		
		\item 
		Confusing \code{=} (assignment) with \code{==} (comparison).
		
		\item 
		Off-by-one errors in loops and arrays. 
		Arrays start at index 0 and ends at $N - 1$.
		
		\item 
		Forgetting that array length is \code{.length} but 
		\code{String} length is \code{.length()}.
		
		\item 
		Null pointer exceptions.
	
		
	\end{itemize}

	\item
	\textbf{Start assignments early.} 
	Use the Help Desk and office hours. 
	Do not just read. Compile and run constantly.
	
	\item
	For exams, practice the sample exam questions under timed conditions. 
	Trace code by hand, predict output, and write small code fragments on paper. 
	If your instructor allows a reference sheet, build it yourself.

\end{itemize}




% ~~~~~~~~~~~~~~~~~~~~~~~~~~~~~~~~~~~~~~
\section{Using AI Responsibly}

Use \AI\ as a tutor, not as a replacement for your own thinking. 
Ask for hints, explanations, and debugging questions. 
\emph{Do not ask \AI\ to write your assignment and then submit it as your own.} 
You should be able to explain every line of code you turn in. 
Follow your instructor’s academic integrity and \AI-use policies.




% ~~~~~~~~~~~~~~~~~~~~~~~~~~~~~~~~~~~~~~
\end{document}